\documentclass[11pt,a4paper]{article}

\usepackage[margin=2.5cm]{geometry}
\usepackage[T1]{fontenc}
\usepackage{lmodern}
\usepackage{amsmath}
\usepackage{booktabs}
\usepackage{tabularx}
\usepackage{enumitem}
\usepackage{xcolor}
\usepackage{listings}
\usepackage{hyperref}

\hypersetup{colorlinks=true, linkcolor=blue!50!black, urlcolor=blue!50!black}
\lstset{basicstyle=\ttfamily\small, columns=fullflexible, breaklines=true, frame=single,
  framerule=0.3pt, rulecolor=\color{black!30}, backgroundcolor=\color{black!3}, showstringspaces=false}
\newcommand{\code}[1]{\texttt{#1}}

\title{CPU--GPU Co-Scheduling Module for CloudSim 7G\\[4pt]
\large Phase 3 Report: Synergy's Scheduling Policies}
\author{Chirag \and Hamiz}
\date{23 November 2026}

\begin{document}
\maketitle

\begin{abstract}
Phase~3 adds Synergy's own scheduling policies: Synergy-Greedy and Synergy-Tune allocation, first-fit
and consolidated placement with SOFT locality, the round-based lifecycle, and the SRTF, LAS and FTF
orderings, together with a \code{synergy-replica} preset that locks every setting to Synergy's simulator.
Before writing code we read Synergy's simulator and found that it differs from the paper in several
places that change results: there is no GPU/CPU/memory sort in Tune, the reference run places jobs by
first fit rather than on the fullest server, progress is counted in whole iterations per round, and
finish times are always round boundaries. We follow the code. A run on a trace built from the paper's
recipe shows the paper's ordering of results (Tune ahead of Greedy, both far ahead of the fixed share)
and takes 22\,s for 6{,}500 jobs and three schedulers. All 261 tests of the project pass.
\end{abstract}

\tableofcontents

%======================================================================
\section{Why read Synergy's code first}
%======================================================================

Phase~4 must reproduce numbers produced by Synergy's simulator (commit \code{1bb2bed}), not numbers
derived from the paper's prose. The full findings, with file and line references, are in
\code{docs/validation-plan.md}. The ones that shaped this phase:

\begin{center}
\begin{tabularx}{\linewidth}{@{}lXX@{}}
\toprule
Topic & Paper says & Code does (and we follow) \\
\midrule
Tune's job order & sort by GPU, then CPU, then memory (\S4.2) & policy order only (FIFO: arrival); a job is
  tried if its GPUs fit (\code{fifo\_synergy\_new.py}) \\
Placement of the reference run & best fit & first fit over free GPUs in index order
  (\code{SYNERGY\_RANDOM}); the placement mode is set by the scheduler's name, not the command line \\
Progress & continuous & \(\lfloor 300 / \text{iteration time} \rfloor\) whole iterations per round \\
Finish time & when work ends & at the end of the round the work ends in \\
Ideal demand & best-case resources & capped at one server, with the speedup re-interpolated \\
Resources & GPU, CPU, memory & also storage speed, but it equals memory in every config and profile \\
\bottomrule
\end{tabularx}
\end{center}

The artifact readme also numbers figures differently from the OSDI paper (its ``Fig~8'' for the
single-GPU FIFO sweep is the paper's Fig.~9); the architecture document follows the readme. The
validation plan lists both numberings.

%======================================================================
\section{What was built}
%======================================================================

\subsection{Throughput model: Synergy's ideal demand}
A job's ideal demand is its GPUs times the profile's ideal cores and memory per GPU. Following
\code{jobs/workload.py}:
\begin{itemize}
  \item if that exceeds one server, it is capped at the server's size and the speedup is
        re-interpolated on the straight line between the fixed share (speedup 1) and the ideal point,
        in cores:
        \[ s' = 1 + \frac{s - 1}{c_\text{ideal} - c_\text{fair}}\,(c_\text{capped} - c_\text{fair}); \]
        for a 4-GPU AlexNet job (ideal 48 cores, fair 12, speedup 2.94) capped at 24 cores, \(s' = 1.647\);
  \item a job with more GPUs than a server gets the fixed share and speedup 1.
\end{itemize}
In \code{synergy\_two\_point} mode a job runs at \(s'\) when it holds at least this demand and at 1
otherwise. Two checks at load time keep this exact: every profile whose ideal demand is at or below the
fixed share has speedup 1 (so a ``+fair'' job never gets a speedup, as in Synergy), and storage demand
equals memory demand.

\subsection{Allocation}
\begin{itemize}
  \item \code{synergy\_greedy}: the ideal demand; a job that does not fit is skipped this round
        (Synergy's scheduler without suffix, paper \S3.3).
  \item \code{synergy\_tune}, a \emph{joint} allocation and placement policy (new plug point
        \code{JointAllocationPlacement}), following \code{cluster.py} \code{\_tune}:
        \begin{enumerate}
          \item skip the job if the cluster has fewer free GPUs than it needs;
          \item place it at its ideal demand;
          \item else, if that is above the fixed share in any dimension, drop it to the fixed share and retry;
          \item else choose servers with enough free GPUs (one server for multi-GPU jobs; otherwise by
                fewest free GPUs) and shrink co-located jobs to the fixed share, trying jobs on fewer
                servers first, then most cores per GPU, then most memory per GPU, until the job's share fits;
          \item place it again; if it still does not fit, Synergy exits, and so does our run.
        \end{enumerate}
        Synergy sorts each server's job list in place, so ties between equally ranked peers depend on
        the order of earlier sorts in the same round; Tune keeps that order per round to match it.
\end{itemize}

\subsection{Placement}
\begin{itemize}
  \item \code{first\_fit} (SOFT by default): Synergy's \code{SYNERGY\_RANDOM}, which is not random:
        servers in index order, taking from each as many GPUs as still fit with their share of cores
        and memory. A 1-GPU job lands on the lowest-numbered server that fits; a multi-GPU job may be
        split.
  \item \code{consolidated} with \code{locality: soft}: Synergy's \code{SYNERGY\_PLACEMENT}. A multi-GPU
        job takes the fullest server that fits it whole; otherwise, and always for a 1-GPU job, it takes
        one GPU at a time from the server with the fewest free GPUs.
  \item Both split a job \emph{in proportion to its GPUs}: each part gets exactly its GPUs' share of
        cores and memory (paper \S4.2), via the shared helper \code{Shares}.
\end{itemize}

\subsection{Lifecycle and Synergy's rounds}
\code{round\_based} re-allocates every job from scratch each round (the plan starts empty); preemption is
free, as in Synergy. Its option \code{synergy\_rounds} reproduces Synergy's time exactly:
\begin{itemize}
  \item each job's speed is quantised to \(\lfloor \text{round} / t \rfloor\) iterations per round, with
        \(t\) = iteration time / penalty / speedup computed in Synergy's order, so the floating-point
        result is identical (e.g. ResNet-18: 468 per round at the fixed share, 1125 at speedup 2.4);
  \item a job whose work ends inside a round is recorded as finishing at that round's end;
  \item jobs whose work is done at a round's start are finished before planning, as Synergy ends leases
        before scheduling.
\end{itemize}
It needs periodic rounds only; the setup refuses anything else. On the mean JCT this changes results by
well under 1\%; it matters for a job-by-job comparison.

\subsection{Orderings}
\begin{center}
\begin{tabularx}{\linewidth}{@{}lX@{}}
\toprule
Ordering & Key (lowest first), as in \code{schedulers/*\_synergy\_new.py} \\
\midrule
\code{srtf} & starved first, remaining time, arrival \\
\code{las} & starved first, attained service in GPU-seconds, arrival \\
\code{ftf} & starved first, (time since arrival + remaining time) / whole run time at the fixed share, arrival \\
\bottomrule
\end{tabularx}
\end{center}
\begin{itemize}
  \item \emph{Starved}: a job not run for 1000 hours gets priority from the next round on (SRTF and LAS
        and FTF; see the Phase~4 report: this was first documented wrongly for FTF).
  \item \emph{Remaining time}: \code{fixed\_share} for ``+fair'' schedulers, \code{current\_or\_ideal} for
        Greedy and Tune. Synergy sorts \emph{before} resetting speedups, so a job that ran at the fixed
        share last round is estimated at speedup 1; \code{current\_or\_ideal} reproduces this.
\end{itemize}

\subsection{The \code{synergy-replica} preset}
\code{presets/synergy-replica/} runs \code{FIFO+tune}, \code{FIFO+fair} and \code{FIFO} on Synergy's
cluster with every setting from Synergy's code (round 300\,s, first fit, two-point throughput, penalty
1, \code{synergy\_rounds}). Its \code{jobs.csv} must be extracted from one Synergy run; this needs all
jobs of the run, including the 266 unfinished when Synergy stopped, which the extractor in
\code{tools/} does not yet keep (an open item for its author).

%======================================================================
\section{Tests}
%======================================================================

\begin{center}
\begin{tabularx}{\linewidth}{@{}lX@{}}
\toprule
Test & Checks \\
\midrule
capped demand & 4-GPU AlexNet: (4, 24, 500) at speedup 1.647; 16-GPU ResNet-18: fixed share, speedup 1 \\
first fit vs fullest & a 1-GPU job goes to server 0 under first fit, to the half-full server 1 under SOFT
  consolidated \\
proportional split & a 6-GPU job over 3 + 3 free GPUs gets (3, 9, 187.5) on each server; STRICT keeps
  it on one server \\
share counting & a server with free GPUs but too few cores is skipped \\
SOFT spread & with 2, 3, 6 free GPUs a 7-GPU job takes 2, 3 and 2, fullest first \\
Greedy & the third ResNet-18 needing 9 cores, with 6 left, is skipped \\
Tune & five 1-GPU ResNet-18 jobs on one server: jobs 1 and 2 get 9 cores, 3 and 4 the fixed share, job 5
  forces job 1 down to the fixed share; only job 2 keeps speedup 2.4 \\
Tune skip & a ninth 1-GPU job on 8 GPUs is skipped \\
SRTF & order flips between fixed-share and ideal estimates, and back when the job ran at the fixed share \\
LAS, FTF & keys as above, worked by hand \\
starvation & a 1001-hour-starved job sorts first from the round after it is marked \\
Synergy rounds & 1000 iterations at 468 per round finish at 900\,s (640\,s continuous); 3000 at 1125
  per round under Tune at 900\,s (800\,s continuous); exactly 936 finish at 600\,s; an arrival at 10\,s
  starts at 300\,s; event-driven rounds are refused \\
preset & synergy-replica's three candidates run end to end; every finish is a multiple of 300\,s; Tune
  beats the fixed share \\
\bottomrule
\end{tabularx}
\end{center}

All existing tests still pass, including the architecture rules (the new helpers \code{Shares},
\code{Starvation} and \code{RemainingEstimate} live in the policy package; no policy implementation uses
another). Project total: 261 tests, 0 failures.

%======================================================================
\section{A first look at scale}
%======================================================================

As a sanity check, not validation, we generated a trace with the paper's recipe (\S5.1): 128 GPUs,
single-GPU jobs, Poisson arrivals at 9 jobs per hour, durations at the fixed share of \(10^x\) minutes
(\(x\) uniform in \([1.5, 3]\) with probability 0.8, else \([3, 4]\)), split (20, 70, 10), 6{,}500 jobs.

\begin{center}
\begin{tabular}{@{}lrr@{}}
\toprule
Candidate & Mean JCT, jobs 4000--4999 (h) & p99 (h) \\
\midrule
FIFO + Tune & 37.1 & 168.7 \\
FIFO + Greedy & 41.0 & 173.1 \\
FIFO + fixed share & 103.5 & 232.0 \\
\bottomrule
\end{tabular}
\end{center}

The ordering matches the paper (Fig.~9: 22\,h Tune, 81\,h fixed share at 9 jobs per hour): Tune is
2.8\(\times\) better than the fixed share here, 3.7\(\times\) in the paper. Absolute values differ because
Synergy's reference workload is not generated from this recipe but built from the Philly trace (a pickle
shipped with its simulator), and at this load the cluster is near saturation, where small changes in load
move JCTs strongly. Replaying Synergy's own job list (Phase~4) removes this difference. The three runs
together took 22\,s (about 11{,}000 rounds each).

%======================================================================
\section{Next step}
%======================================================================
Phase~4: obtain Synergy's job list (one Synergy run plus the extractor option for unfinished jobs), run
\code{synergy-replica}, and compare with the reference: mean JCT within 10\%, and with
\code{synergy\_rounds} a job-by-job match. Then the remaining validation steps, the \code{FIFO+fair}
reference (about 81\,h), and DRF and Tetris.

\end{document}
